\begin{definition}[Normalized integer-valued order]\label{mod:localfield-valuation}
Let $F$ be a nonarchimedean local field, with Mathlib's canonical
\texttt{ValuativeRel} valuation.  Transport that valuation along the local-field
value-group equivalence and take the negative logarithm to define
\[
  \operatorname{ord}_F:F\longrightarrow \mathbf Z\cup\{\infty\}
       =\operatorname{WithTop}(\mathbf Z).
\]
In Lean this function is bundled as an \texttt{AddValuation}.  It satisfies
\[
 \begin{aligned}
  \operatorname{ord}_F(0)&=\infty,&
  \operatorname{ord}_F(1)&=0,\\
  \operatorname{ord}_F(xy)&=\operatorname{ord}_F(x)+\operatorname{ord}_F(y),\\
  \min\{\operatorname{ord}_F(x),\operatorname{ord}_F(y)\}
    &\leq \operatorname{ord}_F(x+y).
 \end{aligned}
\]
Moreover, $\operatorname{ord}_F(x)=\infty$ if and only if $x=0$,
$\operatorname{ord}_F(-x)=\operatorname{ord}_F(x)$, and
\[
 \min\{\operatorname{ord}_F(x),\operatorname{ord}_F(y)\}
   \leq \operatorname{ord}_F(x-y).
\]
If the orders of $x$ and $y$ are distinct, then the order of $x+y$ is
their minimum.  In Lean's totalized field operations and the extended
operations on $\operatorname{WithTop}(\mathbf Z)$,
\[
 \operatorname{ord}_F(x^{-1})=-\operatorname{ord}_F(x),\qquad
 \operatorname{ord}_F(x/y)=\operatorname{ord}_F(x)-\operatorname{ord}_F(y),
\]
and the natural- and integer-power formulas hold.  On nonzero elements these
are the usual local-field identities; at zero they express Lean's conventions
$0^{-1}=0$ and $-\infty=\infty$ in \texttt{WithTop}.  The API also gives
finite-sum lower bounds.
Comparison by $\operatorname{ord}_F$ is the reverse of the canonical valuative relation.
The inequality $0\leq\operatorname{ord}_F(x)$ characterizes membership in the valuation
subring, and, for an element of that subring, $0<\operatorname{ord}_F(x)$ characterizes
membership in its maximal ideal.  Every integer occurs as the finite order of
some element.  Finally, for the canonical discrete valuation,
\[
  \operatorname{ord}_F(\pi)=1
  \quad\Longleftrightarrow\quad
  \pi\text{ is a uniformizer}.
\]

For $n\in\mathbf Z$, define the elementwise depth congruence
\[
  \operatorname{CongruentAtDepth}_n(x,y)
  \quad\Longleftrightarrow\quad
  (n:\operatorname{WithTop}(\mathbf Z))\leq\operatorname{ord}_F(x-y).
\]
At each fixed depth this is an equivalence relation, and it is monotone when
the depth is decreased.  It is preserved by addition, subtraction, negation,
and translation.  If $\operatorname{ord}_F(a)\geq r$, multiplication by $a$ carries depth
$s$ congruence to depth $r+s$ congruence; in particular integral factors
preserve depth.  More precisely, if $a\equiv_n b$, $x\equiv_n y$,
$\operatorname{ord}_F(a)\geq0$, and $\operatorname{ord}_F(y)\geq0$, then
$ax\equiv_n by$.  Inversion preserves depth when both elements have order
zero.  This raw inequality is the order-theoretic form of
the manuscript's congruence modulo $\mathfrak p_F^n$; its identification with
the later fractional-ideal API belongs to the lattice node.
\LeanFile{LanglandsFirstMainLemma/LocalField/Valuation.lean}
\lean{LanglandsFirstMainLemma.ord}
\lean{LanglandsFirstMainLemma.CongruentAtDepth}
\leanok
\SourceText{Local notation, subsection Local fields, ideals, and characters.}
\uses{mod:prelude}
\FormalizationContract
\end{definition}
